\begin{ex}
Mẫu số chung nhỏ nhất của hai phân số $\frac{1}{2}$ và $\frac{1}{3}$ là bao nhiêu?
\shortans{6}
\loigiai{
Mẫu số chung là $2 \times 3 = 6$.
}
\end{ex}

\begin{ex}
Mẫu số chung nhỏ nhất của hai phân số $\frac{1}{4}$ và $\frac{1}{6}$ là bao nhiêu?
\shortans{12}
\loigiai{
12 chia hết cho cả 4 và 6.
}
\end{ex}

\begin{ex}
Mẫu số chung nhỏ nhất của $\frac{3}{5}$ và $\frac{7}{10}$ là bao nhiêu?
\shortans{10}
\loigiai{
10 chia hết cho 5.
}
\end{ex}

\begin{ex}
Quy đồng mẫu số hai phân số $\frac{2}{3}$ và $\frac{3}{4}$ với mẫu số chung 12, phân số $\frac{2}{3}$ trở thành phân số nào?
\shortans{8/12}
\loigiai{
$\frac{2 \times 4}{3 \times 4} = \frac{8}{12}$.
}
\end{ex}

\begin{ex}
Trong bài toán trên, phân số $\frac{3}{4}$ trở thành phân số nào?
\shortans{9/12}
\loigiai{
$\frac{3 \times 3}{4 \times 3} = \frac{9}{12}$.
}
\end{ex}

\begin{ex}
Quy đồng mẫu số $\frac{1}{4}$ và $\frac{5}{8}$ với mẫu số chung 8, phân số $\frac{1}{4}$ trở thành phân số nào?
\shortans{2/8}
\loigiai{
$\frac{1 \times 2}{4 \times 2} = \frac{2}{8}$.
}
\end{ex}

\begin{ex}
Mẫu số chung của $\frac{2}{5}$ và $\frac{3}{7}$ là bao nhiêu?
\shortans{35}
\loigiai{
$5 \times 7 = 35$.
}
\end{ex}

\begin{ex}
Quy đồng $\frac{2}{5}$ với mẫu số chung 35 được tử số mới là bao nhiêu?
\shortans{14}
\loigiai{
$2 \times 7 = 14$.
}
\end{ex}

\begin{ex}
Quy đồng $\frac{3}{7}$ với mẫu số chung 35 được tử số mới là bao nhiêu?
\shortans{15}
\loigiai{
$3 \times 5 = 15$.
}
\end{ex}

\begin{ex}
Mẫu số chung nhỏ nhất của $\frac{5}{6}$ và $\frac{7}{18}$ là bao nhiêu?
\shortans{18}
\loigiai{
18 chia hết cho 6.
}
\end{ex}

\begin{ex}
Quy đồng $\frac{5}{6}$ với mẫu số chung 18 được tử số mới là bao nhiêu?
\shortans{15}
\loigiai{
$5 \times 3 = 15$.
}
\end{ex}

\begin{ex}
Mẫu số chung nhỏ nhất của $\frac{1}{2}, \frac{1}{3}$ và $\frac{1}{4}$ là bao nhiêu?
\shortans{12}
\loigiai{
12 chia hết cho 2, 3, 4.
}
\end{ex}

\begin{ex}
Quy đồng $\frac{1}{2}$ với mẫu số chung 12 được tử số là bao nhiêu?
\shortans{6}
\loigiai{
$1 \times 6 = 6$.
}
\end{ex}

\begin{ex}
Quy đồng $\frac{1}{3}$ với mẫu số chung 12 được tử số là bao nhiêu?
\shortans{4}
\loigiai{
$1 \times 4 = 4$.
}
\end{ex}

\begin{ex}
Quy đồng $\frac{1}{4}$ với mẫu số chung 12 được tử số là bao nhiêu?
\shortans{3}
\loigiai{
$1 \times 3 = 3$.
}
\end{ex}

\begin{ex}
Mẫu số chung nhỏ nhất của $\frac{3}{8}$ và $\frac{5}{12}$ là bao nhiêu?
\shortans{24}
\loigiai{
24 chia hết cho 8 và 12.
}
\end{ex}

\begin{ex}
Tử số mới của $\frac{3}{8}$ khi quy đồng mẫu số chung 24 là bao nhiêu?
\shortans{9}
\loigiai{
$3 \times 3 = 9$.
}
\end{ex}

\begin{ex}
Tử số mới của $\frac{5}{12}$ khi quy đồng mẫu số chung 24 là bao nhiêu?
\shortans{10}
\loigiai{
$5 \times 2 = 10$.
}
\end{ex}

\begin{ex}
Quy đồng mẫu số $\frac{4}{9}$ và $\frac{1}{3}$ với mẫu số chung 9, phân số $\frac{1}{3}$ trở thành phân số nào?
\shortans{3/9}
\loigiai{
$\frac{3}{9}$.
}
\end{ex}

\begin{ex}
Mẫu số chung của $\frac{7}{15}$ và $\frac{3}{5}$ là bao nhiêu?
\shortans{15}
\loigiai{
15.
}
\end{ex}

\begin{ex}
Tử số mới của $\frac{3}{5}$ khi quy đồng với mẫu số chung 15 là bao nhiêu?
\shortans{9}
\loigiai{
$3 \times 3 = 9$.
}
\end{ex}

\begin{ex}
Mẫu số chung nhỏ nhất của $\frac{5}{14}$ và $\frac{3}{7}$ là bao nhiêu?
\shortans{14}
\loigiai{
14.
}
\end{ex}

\begin{ex}
Mẫu số chung nhỏ nhất của $\frac{1}{6}$ và $\frac{1}{9}$ là bao nhiêu?
\shortans{18}
\loigiai{
18 chia hết cho cả 6 và 9.
}
\end{ex}

\begin{ex}
Tử số mới của $\frac{1}{6}$ khi mẫu số chung là 18 là bao nhiêu?
\shortans{3}
\loigiai{
$1 \times 3 = 3$.
}
\end{ex}

\begin{ex}
Tử số mới của $\frac{1}{9}$ khi mẫu số chung là 18 là bao nhiêu?
\shortans{2}
\loigiai{
$1 \times 2 = 2$.
}
\end{ex}
